% ============================================================
% MODULO 08 - CONCLUSAO
% ============================================================
\section{Conclus\~ao}

Este trabalho sustenta uma tese deliberadamente modesta: em pr\'e-escolares com TDAH misto, a ecologia do brincar -- onde, com quem e por quanto tempo ao ar livre -- merece lugar na formula\c{c}\~ao cl\'inica, ao lado dos sintomas, da fam\'ilia e da escola \cite{bronfenbrenner1996}. Uma ``dose verde-social'' complementar -- caminho arborizado para a escola, quintal, recreio em gramado, janela com vista verde para as tarefas -- custa pouco, n\~ao tem efeitos adversos conhecidos e \'e coerente com a literatura \cite{taylor2011,hood2024,fabertaylor2009}, desde que n\~ao substitua o acompanhamento j\'a em curso e n\~ao seja confundida com indica\c{c}\~ao medicamentosa, que este texto n\~ao faz.

Como desdobramentos pr\'aticos para a cl\'inica e a orienta\c{c}\~ao de fam\'ilias, sugerem-se quatro cuidados: (a) registrar a dose di\'aria de brincar externo social (minutos, companhia, local), transformando o registro em instrumento de conversa nas sess\~oes; (b) priorizar espa\c{c}os verdes e abertos quando h\'a predom\'inio de hiperatividade, em sintonia com a literatura; (c) envolver a escola no mapeamento do recreio e do entorno; e (d) manter postura n\~ao diretiva, deixando a pr\'opria crian\c{c}a apontar o que o brincar mobiliza nela.

As limita\c{c}\~oes deste relato s\~ao conhecidas: caso \'unico, aus\^encia de grupo controle e de escala padronizada, dose de brincar e de telas n\~ao medida, informa\c{c}\~ao sobre medica\c{c}\~ao e sono indispon\'ivel e poss\'ivel vi\'es no relato materno. A agenda de continuidade inclui: (1) aplicar SNAP-IV com pais e professores agora e em tr\^es meses; (2) manter por duas semanas um di\'ario simples (hora de rua, com quem, onde, hora de tela, hora de dormir); (3) conduzir de 8 a 12 sess\~oes com protocolo axlineano registrado em vinhetas; (4) realizar entrevista semiestruturada com a m\~ae; e (5) se a pesquisa for institucional, submet\^e-la a Comit\^e de \'Etica em Pesquisa com protocolo registrado.

A contribui\c{c}\~ao, ainda que pequena, \'e genu\'ina: documentar em portugu\^es, sob \'otica humanista, um fen\^omeno cotidiano quase ausente da literatura -- o brincar de vizinhan\c{c}a n\~ao estruturado aos 5 anos -- e oferecer a outros cl\'inicos categorias para escutar, nas sess\~oes, mudan\c{c}as que as escalas ainda n\~ao medem.
